% Stage: — ／ 説明用コード: —
\chapter{エラー訂正フロー}\label{ch:qec_flow}

\begin{chaptergoal}
  \begin{itemize}
  \item シンドローム測定を時間軸上で繰り返す「符号サイクル」と、$T_1$ 緩和時間による制約を理解します。
  \item 物理的な即時訂正を避け、古典コンピュータ上でエラーを追跡する「パウリフレーム」の考え方を理解します。
  \item クリフォード操作がパウリフレーム追跡と相性がよい理由と、非クリフォード操作がその前提を崩す理由を理解します。
  \item 魔法状態テレポーテーションに伴うデコード遅延とストール問題を理解します。
  \item 層を横断する論理エラー率を評価するために、統合シミュレータが必要になる理由と本書全体の狙いを理解します。
  \end{itemize}
\end{chaptergoal}

第3章では、シンドローム測定を1回きりの操作として、いわば静止画のように取り上げました。
第4章では、その静止画の中身であるスタビライザ形式・符号距離・CSS条件・表面符号の幾何を確認しました。
本章では視点を変え、実際にFTQCで計算を実行する際にこの静止画がどのように\textbf{時間軸上で繰り返され}、論理演算とどう絡み合うのかを見ていきます。
符号サイクルの繰り返しと $T_1$ 制約、物理的な訂正を行わない代わりに古典ソフトウェアでエラーを追跡する仕組み（パウリフレーム）、クリフォード操作と非クリフォード操作の非対称性、そして非クリフォード操作にまつわるデコード遅延とストール問題を順に確認し、最後にこれらすべてを一つの評価軸でまとめるために本書がなぜハイブリッド型シミュレータを組み立てるのかを述べて、第I部を締めます。

\section{誤り訂正サイクルと \texorpdfstring{$T_1$}{T1} 制約}\label{sec:cycle_t1}

\subsection{符号サイクル}

第3章・第4章のシンドローム測定は、1回の測定でどれだけの情報が得られるかという説明に専念していました。
表面符号などの能動的な量子メモリでは、計算中も保存中もこの測定を周期的に繰り返します。
1回のシンドローム測定（パリティ検査）を行ってから次のシンドローム測定を行うまでの一続きの量子回路を、\textbf{符号サイクル}と呼びます。

論理量子ビットに何らかの論理演算を実行させている間も、演算をせずにただ情報を保持させている間（アイドリング）も、符号サイクルは休みなく繰り返し実行されます。
各サイクルの量子回路の内容は、実行中の論理演算の種類に応じて多少変化しますが、どの場合であっても、量子ビットが物理的に保持できる時間に対して十分短い時間で完了できる回路深さでなければならない、という制約は共通しています。

\subsection{$T_1$ 緩和時間との競合}

物理量子ビットが励起状態のエネルギーを失う振幅減衰の時定数を\textbf{$T_1$ 緩和時間}、位相コヒーレンスが失われる時定数を\textbf{$T_2$} と呼びます。Markov 的な単純模型では励起状態の生存確率は $\exp(-t/T_1)$ のように連続的に減衰し、$T_1$ を境に突然崩壊するわけではありません\cite{krantz2019guide}。

符号サイクルの役割は、この崩壊が訂正不可能なほど進行する前に誤りを検知することです。
振幅減衰だけを見た設計上の目安として、符号サイクル1回にかかる時間 $t_c$ は
\begin{equation}\label{eq:cycle_vs_t1}
  t_c \ll T_1
\end{equation}
となることが望まれます。ただしこれは閾値条件そのものではなく、$T_2$、漏れ、ゲート・測定誤り、相関誤りを含む1サイクル当たりの実効ノイズが、採用する符号とデコーダの閾値を下回ることが本質です\cite{krantz2019guide,dennis2002topological}。

もし符号サイクルの間隔を延ばしたり、サイクルそのものを休止したりすると、その休止時間中にも物理誤りは確率的に発生し続けます。
独立な Markov 誤りを仮定し、休止時間が各コヒーレンス時間より十分短い領域では、期待故障数は時間にほぼ比例して増えます。長時間や相関ノイズではこの線形近似は成り立たず、故障数が訂正能力を超える確率を評価する必要があります。
これが、能動的な表面符号メモリでは\textbf{論理演算を行っていないアイドリング中も符号サイクルを継続する}理由です。
データ量子ビットは、常にシンドローム測定という「保守作業」を受け続けることによって、はじめて論理情報を長時間保持できます。

\begin{notebox}
  古典のDRAMメモリも、電荷が漏れて情報が消えるのを防ぐために一定間隔で\textbf{リフレッシュ}を行います。
  符号サイクルによるアイドリングの保護は、この古典的なリフレッシュの量子版だと考えると直感的です。
  ただし量子の場合、リフレッシュのたびに情報を直接読み出す（測定する）わけではなく、第3章で見た間接測定＝シンドローム測定を使う点が異なります。
\end{notebox}

\section{パウリフレームと古典ソフトウェア追跡}\label{sec:pauli_frame}

\subsection{物理的な訂正をしないという選択}

符号サイクルでシンドロームが得られ、デコーダが観測結果と整合する訂正の同値類を選んだとしても、素朴には「推定した誤りを打ち消すパウリ操作を物理量子ビットに実際にかける」ことで訂正が完了すると考えたくなります。シンドロームだけから実際に起きた物理誤りを一意に知ることはできません。
しかし、この\textbf{物理的な訂正}には見過ごせないコストがあります。

\begin{itemize}
  \item 訂正操作そのものが、新たな物理誤りの追加源になり得る。
  \item デコードには一定の時間がかかるため（5.4節のストール問題）、訂正操作の完了を待つ間、次の符号サイクルの開始が遅れる。
  \item 論理演算のたびに関係する全量子ビットへ訂正をかけ直すのは非効率である。
\end{itemize}

FTQCの実用的なアーキテクチャでは、多くのパウリ補正を直ちに物理適用せず、後続ゲートと測定の解釈へ繰り延べられます。
代わりに、デコーダが選んだ訂正を適用したものとして解釈するためのパウリ演算子を\textbf{古典コンピュータ上の帳簿}として記録し続け、必要になった時点（測定結果の解釈や論理読み出しの直前など）で帳尻を合わせます。
この帳簿を\textbf{パウリフレーム}と呼びます。

\subsection{パウリフレームが成立する根拠}

パウリフレームが機能する根拠は、クリフォード操作が持つ次の性質です。
クリフォード操作 $U_C$ と任意のパウリ演算子 $P$ に対して
\begin{equation}\label{eq:clifford_conj}
  U_C P = P' U_C, \qquad P' = U_C P U_C^\dagger
\end{equation}
が成り立ち、$P'$ は符号（$\pm1$ の位相）を除いて再びパウリ演算子になります。
これは「$U_C$ がパウリ群を自分自身へ写す」という性質そのもので、クリフォード操作の定義です\cite{aaronson2004stabilizer}。

いま、本来の状態 $\ket\psi$ にパウリ誤り $E$ が乗った状態 $E\ket\psi$ にクリフォード操作 $U_C$ を適用したとします。式\eqref{eq:clifford_conj}より
\begin{equation}
  U_C \bigl(E\ket\psi\bigr)
  = \bigl(U_C E U_C^\dagger\bigr) U_C \ket\psi
  = E' \, U_C\ket\psi, \qquad E' = U_C E U_C^\dagger
\end{equation}
となります。右辺は「誤りのない状態 $\ket\psi$ に $U_C$ をかけたあと、別のパウリ誤り $E'$ が乗っている」と読めます。つまり、

\begin{quote}
  「誤りが乗った状態に演算をかけた結果」と「誤りのない状態に演算をかけたあと、演算に応じて誤りの種類を古典的に付け替えた結果」は、量子状態として完全に等価である
\end{quote}

ということです。この付け替えは行列の積を計算するだけの古典計算であり、実際の量子ビットには一切手を触れません。

この関係を毎サイクル繰り返し適用すれば、量子デバイスは物理補正を毎回実行せずに演算を続けられ、古典コンピュータの側だけがデコーダの選んだ訂正代表を表すパウリ演算子 $F$（各量子ビットに $I,X,Y,Z$ のいずれかを割り当てたもの）を式\eqref{eq:clifford_conj}に従って更新できます。この $F$ が\textbf{パウリフレーム}です。
シンドローム測定の測定値（レコード）についても、その値が本来ゼロになるはずかどうかを、実際の測定値とパウリフレーム $F$ とを比較して古典的に判定します。

\begin{notebox}
  パウリフレームは、量子ビットの「実体」を書き換える代わりに、実体と設計図のあいだの\textbf{ずれ}だけを古典コンピュータに覚えさせておく仕組みだと考えると直感的です。
  設計図どおりに動いていないと分かった時点でいちいち実体を直しに行くのではなく、「今どれだけずれているか」を台帳に記録しておき、最終的な読み出しや論理演算の解釈のときだけその台帳を参照して辻褄を合わせます。
\end{notebox}

\section{論理ゲートの分類とクリフォードの特異性}\label{sec:gate_classification}

\subsection{クリフォードと非クリフォード}

論理演算は、パウリフレームとの相性という観点から大きく2種類に分類されます。

\begin{itemize}
  \item \textbf{クリフォード操作}（$H$、$S$、CNOT、およびこれらの合成）：式\eqref{eq:clifford_conj}のとおり、パウリ演算子をパウリ演算子（符号を除く）に写します。
  \item \textbf{非クリフォード操作}（$T$ ゲートなど）：この写像がパウリ群の外に出てしまい、単純な古典的付け替えでは追跡できません。
\end{itemize}

シンドローム測定の回路自体も、スタビライザ状態の初期化、Clifford ゲート、Pauli 測定という\textbf{スタビライザ操作}で構成されています。初期化と測定はユニタリな Clifford ゲートではありません。
したがって、符号サイクルを回している間に生じるパウリ誤りは、常に式\eqref{eq:clifford_conj}の関係でパウリフレームへ吸収できます。
これが、誤り訂正サイクルとクリフォードな論理演算を、量子デバイスを止めることなく古典ソフトウェアだけで整合させられる理由です。

\subsection{$T$ ゲートがパウリ群を保たないことの確認}

非クリフォードの代表例である $T=\operatorname{diag}(1,\ee^{\ii\pi/4})$ が、なぜパウリ演算子をパウリ演算子に写さないのかを具体的に確かめます。
$TXT^\dagger$ を計算すると
\begin{equation}\label{eq:t_not_clifford}
  TXT^\dagger
  = \begin{pmatrix} 0 & \ee^{-\ii\pi/4} \\ \ee^{\ii\pi/4} & 0 \end{pmatrix}
  = \cos\frac{\pi}{4}\, X + \sin\frac{\pi}{4}\, Y
  = \frac{1}{\sqrt2}(X+Y)
\end{equation}
となります。右辺はパウリ演算子 $X$ と $Y$ の実数係数の重ね合わせであり、単一のパウリ演算子ではありません。
したがって $T$ は式\eqref{eq:clifford_conj}の意味でのクリフォード条件を満たさず、$T$ ゲートの前後でパウリフレームを単純な文字列の付け替えとして押し出すことができません。

このことは、より根本的な事実、すなわちトランスバーサル操作だけではユニバーサルな計算ができないという定理（第4章で触れた話題と表裏一体）とも整合します。
クリフォード操作だけでは複雑なアルゴリズムを実行できないため非クリフォード操作は必須ですが、式\eqref{eq:t_not_clifford}が示すとおり、それを安全に実行するにはパウリフレームの付け替えとは別の仕掛けが必要になります。次節でこれを見ます。

\section{非クリフォード操作のデコード遅延とストール問題}\label{sec:stall}

\subsection{ゲートテレポーテーションとフィードフォワード}

論理 $T$ ゲートは、物理量子ビットに直接かけることができません。
代わりに、あらかじめ $T\ket{+}$ に相当する\textbf{魔法状態}を論理レベルで用意し、対象の論理量子ビットと量子もつれさせたうえで対象を測定する\textbf{ゲートテレポーテーション}という手法で実行します\cite{bravyi2005magic, gidney2019factory}（符号化・蒸留の具体的な手順は第III部で扱います）。

この手順を論理レベルの回路として書くと図\ref{fig:t_teleport}のようになります。

\begin{figure}[htbp]
  \centering
  \begin{quantikz}
    \lstick{$\ket{\psi}_L$}          & \ctrl{1} & \gate{S^{m}} & \rstick{$T\ket{\psi}_L$} \\
    \lstick{$\ket{A}_L=T\ket{+}_L$}  & \targ{}  & \meter{}     &
  \end{quantikz}
  \caption{$T$ ゲートテレポーテーションの概念図（論理レベル）。データ側を制御、魔法状態 $\ket{A}_L$ を標的として $\mathrm{CNOT}$ をかけたのち、\textbf{魔法状態側を $Z$ 基底で測定}します。測定結果 $m\in\{0,1\}$ に応じて、\textbf{データ側}に補正 $S^m$ をかけるかどうかが決まります（フィードフォワード）。目的の状態はデータ側の線に得られます。}
  \label{fig:t_teleport}
\end{figure}

同図のとおり、測定結果 $m$ に応じて後続に補正操作 $S^m$ をかけるかどうかが変わります。
この「測定結果を見てから後続の回路の内容を選ぶ」処理を\textbf{フィードフォワード}と呼びます。

測定を含む Clifford 実装にも結果依存の分岐はあり得ますが、結果依存の\emph{Pauli}補正なら通常はパウリフレームへ吸収できます。一方 $S^m$ は Pauli ではないため、通常のパウリフレームには追加できません。物理的または適応的に実行するか、より一般の\textbf{Clifford frame}として追跡します\cite{chamberland2018frame}。
しかし、$m$ の値を決める測定そのものにも物理誤りが乗っており、真の $m$ を確定するには、その符号サイクルまでの\textbf{シンドローム測定の結果をすべて集めてデコードを完了させる}必要があります。
続けて別の非クリフォード操作が来る場合、どちら向きの補正が要るかを誤ると、以降のパウリフレームの管理そのものが破綻してしまいます。
そのため、後続操作の選択がこの値に依存する地点では、\textbf{必要なデコードを完了させて確定した値を使う}場面が生じます。Clifford frame、選択的テレポーテーション、バッファリングなどで依存関係を先送りできる設計もあり、すべての非 Clifford 境界が必ず停止点になるわけではありません\cite{chamberland2018frame}。

\subsection{リアクション時間とストール}

制御装置が論理パウリ測定（あるいは非クリフォード操作の一部としての測定）を送出してから、デコードを経て確定した測定値が返るまでの時間を\textbf{リアクション時間}と呼びます。
符号サイクル自体は1回あたり短時間（5.1節の $t_c$）で完了しますが、シンドロームから誤りを推定するデコード計算は一般に計算量的に重い問題であり、符号サイクルよりずっと長くかかることがあります。
実際のリアクション時間は、符号距離、デコーダ、古典ハードウェア、制御系の通信経路に依存します。

このリアクション時間の間、量子デバイス側は次の非クリフォード操作へ進めず待たされます。この待機を\textbf{ストール}と呼びます。
ストールしている間も、待機中の論理量子ビットに対する符号サイクルは（5.1節のとおり）回り続けますが、もしストール時間が $T_1$ に対して無視できないほど長くなれば、待機中の量子ビットに蓄積する実効的な誤り率が上昇し、論理エラー率の悪化につながります。
つまりストール問題は、「符号サイクルは $T_1$ より十分短くなければならない（式\eqref{eq:cycle_vs_t1}）」という制約と、「非クリフォードの境界ではデコードの完了を待たなければならない」という制約が正面から衝突する点に生じる、FTQCアーキテクチャ特有のボトルネックです。

\subsection{3つのボトルネックとしての整理}

以上の議論は、FTQC全体の実行時間を大まかに次の3種類のボトルネックのどれかとして捉える見方につながります。

\begin{itemize}
  \item \textbf{魔法状態の供給}：必要な魔法状態数 $N_M$ と生成周期 $T_M$ の積 $N_M T_M$。
  \item \textbf{リアクションの連鎖}：非クリフォード操作どうしの依存関係の長さ $N_R$ とリアクション時間 $T_R$ の積 $N_R T_R$。
  \item \textbf{符号サイクルの繰り返し}：計算全体で必要な符号サイクル数 $N_C$ と、サイクル時間 $T_C$ の積 $N_C T_C$。
\end{itemize}

このうちどれか1つが極端に支配的であれば、計算機の設計とアプリケーションの要求がかみ合っていないことを意味します。
ボトルネックの所在を特定し、デコーダの高速化・コンパイル時の非クリフォード演算の並列化・魔法状態工場の規模調整などで手当てするのが、FTQCシステム設計の実務的な課題です。
本書のシミュレータ（第III部以降）でも、この3つの視点は論理エラー率・実行時間の評価軸として繰り返し登場します。

\section{統合評価環境の必要性と本書の狙い}\label{sec:integrated_sim}

\subsection{層を横断する論理エラー率}

ここまでで確認したFTQCの実行フローは、少なくとも次の3層が互いに絡み合って動いています。

\begin{itemize}
  \item \textbf{量子誤り訂正層}：符号サイクルによるシンドローム測定と、そこからのデコード（第II部以降）。
  \item \textbf{クリフォード論理演算層}：パウリフレームによる古典的な追跡（本章5.2–5.3節）。
  \item \textbf{非クリフォード層}：魔法状態を介したゲートテレポーテーションとフィードフォワード（本章5.4節）。
\end{itemize}

最終的にアプリケーションが気にする量は、これら3層すべてを通過したあとの\textbf{論理エラー率}です。
ところが、既存の評価環境は層ごとに分断されています。
たとえば第6章で導入する Stim はスタビライザ回路の高速シミュレーションに特化した強力なツールですが、扱えるのはクリフォード演算に限られたクリフォードシミュレータです。
非クリフォード操作を含む厳密な量子状態を追いたければ状態ベクトルシミュレータが必要になりますが、これは量子ビット数に対して指数的なコストがかかり、量子誤り訂正符号が要求する物理量子ビット数をそのまま扱うことはできません。
つまり、単独のツールでは3層をまとめて評価できないというのが現状です。

\subsection{Stim + 状態ベクトル + パウリフレームのハイブリッド}

本書で組み立てるシミュレータは、この3層を次のように役割分担させることで、それぞれの限界を回避します。

\begin{itemize}
  \item \textbf{Stim}：スタビライザ回路を高速にサンプリングし、検出器誤りモデルを生成します。復号そのものは PyMatching や BP+OSD など外部デコーダが担います\cite{gidney2021stim}。
  \item \textbf{Qulacs（状態ベクトル）}：魔法状態や非 Clifford 操作が関わる、実装が選んだ小規模な論理状態を状態ベクトルとして追跡します。
  \item \textbf{パウリフレーム（ソフトウェア追跡）}：本章で見た古典的な帳簿を実際にプログラム上で実装し、上の2つのバックエンドの結果を、実際に物理操作をかけ直すことなく1つの整合した論理エラー率へまとめ上げます。
\end{itemize}

つまり「大部分をクリフォードシミュレータで高速に流し、非クリフォードが絡む狭い範囲だけを状態ベクトルで厳密に扱い、両者のつなぎ目をパウリフレームで古典的に処理する」というのが、本書のハイブリッド型アーキテクチャの骨格です。
第I部で確認した4つの制約（第2章）、シンドローム測定の必然性（第3章）、符号の幾何（第4章）、そして本章で見たフロー全体は、いずれもこのハイブリッド構成がなぜ必要かを説明するための土台でした。

\subsection{読書の地図：Stage 01 から F へ}

本書はこのシミュレータを一気に完成形として提示するのではなく、次の表のように段階（Stage）を踏んで組み立てます。

\begin{center}
\begin{tabular}{cll}
\hline
Stage & 到達点 & 対応する部 \\
\hline
基礎 & FTQC と誤り訂正の考え方 & 第I部 \\
01 & 現象論 MPP による論理演算の骨格 & 第II・III部 \\
02 & 回路レベル格子手術 & 第IV部 \\
03 & $d=3$ 固定の統合シミュレータ & 第V部 \\
\hline
\end{tabular}
\end{center}

第II部ではまず、この骨格の入口として Stim の基本的な使い方（第6章）、表面符号パッチの作り方（第7章）、デコードと評価基準の考え方（第8章）を押さえます。
そこから第III部以降で、本章で述べたパウリフレーム・クリフォード/非クリフォードの区別・ストール対策の考え方を、実際に動くコードとして段階的に積み上げていきます。
本章では実装の詳細には踏み込まず、以降の章でコードとして再会するための見取り図を示すことに専念しました。

\subsection{統合評価環境の 3 つの目的と本書の到達範囲}\label{sec:integrated_goals}

層を横断する統合評価環境には、個別のツールでは見えない問題を調べる役割があります。
目的は次の 3 つに整理できます。
一つ目は、$T$ ゲートのたびに生じる「測定 $\to$ 古典側への転送 $\to$ デコード $\to$ フィードフォワード」の往復時間が、量子ビットの寿命に対してどれだけ余裕を持つかを定量化することです。
二つ目は、論理的な計算時間だけでなく、魔法状態の生成時間やストールによる待ちを含めた実効的な計算時間を評価し、デコーダの高速化・通信帯域・工場の並列度のどこに資源を投じるべきかを判断することです。
三つ目は、測定結果によって回路が変わる環境で、限られた接続性のもとでの演算の配置と経路を決めるコンパイラを試すテストベッドになることです。

本書のシミュレータがこの 3 つのどこまでを扱い、どこから先を扱わないかを表\ref{tab:ch5_integrated_goals}に示します。
本書は実時間を持たない論理層のシミュレータなので、時間に関わる量はラウンド数と解析的なコストモデルで扱い、実機の遅延そのものは測りません。

\begin{table}[htbp]
  \centering
  \caption{本書が扱う評価軸と対象外の実測量}
  \label{tab:ch5_integrated_goals}
  \begin{tabularx}{\textwidth}{>{\raggedright\arraybackslash}p{0.25\textwidth}X X}
    \hline
    \textbf{評価軸} & \textbf{本書で扱うもの} & \textbf{本書で扱わないもの} \\
    \hline
    論理誤り & 検出器、復号結果、論理残差 & 実機での論理誤り率の保証 \\
    回路コスト & ラウンド数、補助量子ビット、簡易モデル & 実時間のデコード・通信遅延 \\
    非 Clifford 操作 & 魔法状態、ゲートテレポーテーション、QPD & 実機の工場運用・資源配分 \\
    \hline
  \end{tabularx}
\end{table}

\section{まとめ}\label{ch5:summary}

\begin{itemize}
  \item シンドローム測定を1回きりの静止画としてでなく、\textbf{符号サイクル}として時間軸上に繰り返す視点が必要です。$t_c\ll T_1$ は有用な設計目安ですが、実効ノイズには $T_2$、ゲート・測定誤り、漏れ、相関も含まれます。
  \item FTQCは物理的な即時訂正を避け、代わりに\textbf{パウリフレーム}という古典的な帳簿でエラーを追跡します。この仕組みは、クリフォード操作がパウリ演算子をパウリ演算子に写す性質（式\eqref{eq:clifford_conj}）に支えられています。
  \item 非クリフォード操作（$T$ ゲート）はこの性質を持たず（式\eqref{eq:t_not_clifford}）、魔法状態を介したゲートテレポーテーションとフィードフォワードによって実行されます。
  \item デコードにかかる時間（リアクション時間）が符号サイクルより長いため、非クリフォード操作の境界で計算がストールする可能性があり、これがFTQCの主要なボトルネックのひとつです。
  \item 量子誤り訂正・クリフォード論理演算・非クリフォード層を横断して論理エラー率を評価するには、Stim・状態ベクトル・パウリフレームを組み合わせたハイブリッドシミュレータが必要であり、これが本書全体の狙いです。
\end{itemize}

第I部はここで一区切りです。次章からは第II部に入り、実際に手を動かしてStimの基本的な使い方から本書のシミュレータの構築を開始します。

\begin{exercise}
  \begin{enumerate}
    \item $t_c\ll T_1$ が有用な設計目安である理由と、それだけでは閾値以下を保証しない理由を、アイドリング中の $T_2$、ゲート・測定誤りにも触れて説明しなさい。
    \item クリフォード操作 $U_C$ とパウリ演算子 $E$ について、$U_C(E\ket\psi) = E' U_C \ket\psi$（$E'=U_CEU_C^\dagger$）が成り立つことが、パウリフレーム追跡にとってなぜ本質的なのかを一文で説明しなさい。
    \item 式\eqref{eq:t_not_clifford}の計算を確認し、$T$ ゲートが式\eqref{eq:clifford_conj}の意味でクリフォードでないことを、$TZT^\dagger$ についても同様に確かめなさい。
    \item リアクション時間がストールを引き起こす条件を、符号サイクル時間・$T_1$・デコード時間の関係として説明し、ストールを緩和する方法を1つ挙げなさい。
  \end{enumerate}
\end{exercise}
